%! The Statistical Cycle and Types of Data — Unit 1, Lesson 1.1 — Slides (Beamer)
% PILOT SHAPE (2026-09-06), Main Ideas / Notes table (2026-09-12).
% GRADUAL RELEASE deck: title → targets → warm-up → hook (unresolved) → I DO divider → two or
% three frames per notes ROW, labelled by the row's name and problem numbers → WE DO (a live
% surface, un-answered) → spoken debrief → close & START THE HOMEWORK, which is the individual
% practice. No group activity, no practice launch, no exit ticket. There is no spoiler rule:
% the notes frames ARE the direct instruction.
\documentclass[aspectratio=169, 11pt]{beamer}
\usepackage{saar-beamer}   % provides \ceruleanheader{} and \sectionlabel[color]{}

\newcommand{\redink}[1]{{\color{redacc}\bfseries #1}}

\begin{document}

% TITLE SLIDE (CourseName is not defined in beamer, so the name is literal)
{
\setbeamercolor{background canvas}{bg=cerulean}
\begin{frame}[plain]
  \vspace{2.8cm}\hspace{0.55cm}%
  \begin{minipage}{0.9\textwidth}
    {\color{goldacc}\bfseries\small LESSON 1.1}\\[0.5cm]
    {\color{white}\bfseries\LARGE The Statistical Cycle and Types of Data}\\[1.2cm]
    {\color{frostmid}\small Statistical Analysis \& Algebraic Reasoning~$\cdot$~Unit 1}
  \end{minipage}
\end{frame}
}

% ── LEARNING TARGETS ─────────────────────────────────────────────────────────
\begin{frame}[t, plain]
  \ceruleanheader{Today's targets}
  \sectionlabel{Lesson 1.1}
  By the end of today, I can\ldots
  \begin{itemize}\setlength{\itemsep}{5pt}
    \item name the four stages of the \textbf{data cycle} and say what each hands to the next.
    \item test a question against three checks --- answers \textbf{vary}, the \textbf{group}
          is named, the \textbf{measurement} is named --- and repair one that fails.
    \item decide from the \emph{question alone} whether the data will be
          \textbf{categorical} or \textbf{quantitative}, and build a \textbf{frequency table}.
    \item write a \textbf{conclusion in context} --- and name the \textbf{stage} a flawed
          study broke at.
  \end{itemize}
  \begin{block}{How today works}
    Warm-up $\to$ \textbf{notes together, row by row} $\to$ \textbf{one survey we work as a
    class} $\to$ debrief $\to$ \textbf{you start the homework here, alone}.
  \end{block}
\end{frame}

% ── WARM-UP ──────────────────────────────────────────────────────────────────
\begin{frame}[t, plain]
  \ceruleanheader{Warm-Up --- 5 minutes, silent}
  \sectionlabel{the Student Council's survey}
  \vspace{0.3cm}
  \begin{center}
    \begin{minipage}{0.86\textwidth}
      Last spring the council asked \textbf{60} Riverbend students how they usually get to
      school. \textbf{27} said the bus. Riverbend has \textbf{900} students.
      \begin{enumerate}\setlength{\itemsep}{7pt}
        \item What percent of the students asked ride the bus?
        \item Predict how many of the 900 ride the bus.
        \item Classify two variables --- \emph{how the student gets to school}, and
              \emph{minutes the trip takes} --- and say how you know.
      \end{enumerate}
    \end{minipage}
  \end{center}
  \vspace{0.4cm}
  \begin{center}
    {\color{cerulean}\bfseries Keep the page --- the notes open on these same 60 students.}
  \end{center}
\end{frame}

% ── HOOK — leave it UNRESOLVED; problem 14 (row 4) settles it ────────────────
{
\setbeamercolor{background canvas}{bg=cerulean}
\begin{frame}[plain]
  \vspace{1.2cm}
  \begin{center}
    {\color{frostmid}\bfseries\small THE STUDENT COUNCIL WANTS MORE BIKE RACKS.}\\[0.7cm]
    {\color{white}\bfseries\Large 60 students asked how they get to school. \\[0.3cm]
     6 said bike. ``10\% bike now --- so a lot more \emph{would} if there were racks.''}\\[0.8cm]
    {\color{goldacc}\bfseries\large The arithmetic is right. Does the survey support the claim?}\\[0.4cm]
    {\color{frostmid}\small Vote: yes, no, or can't tell. We settle this at problem 14.}
  \end{center}
\end{frame}
}

% ── I DO — direct instruction begins ─────────────────────────────────────────
{
\setbeamercolor{background canvas}{bg=cerulean}
\begin{frame}[plain]
  \vspace{1.8cm}\hspace{0.8cm}%
  \begin{minipage}{0.86\textwidth}
    {\color{goldacc}\bfseries\small GUIDED NOTES \ \textbullet\ \ I DO}\\[0.7cm]
    {\color{white}\bfseries\Large Open to your Guided Notes page. We work the table row by
    row --- fill each term as we name it, not before.}\\[0.9cm]
    {\color{frostmid}\small Four terms go in the vocabulary box today. Write each one the
    moment we use it.}
  \end{minipage}
\end{frame}
}

% ── ROW 1 — THE DATA CYCLE ───────────────────────────────────────────────────
\begin{frame}[t, plain]
  \ceruleanheader{Four stages, in order --- and then it starts over}
  \sectionlabel{notes row 1 $\cdot$ Data Cycle $\cdot$ problems 1--4}
  \vspace{2pt}
  \begin{center}
    \small
    \renewcommand{\arraystretch}{1.3}
    \begin{tabular}{@{} c l l @{}}
    \hline
    \textbf{Stage} & \textbf{What happens} & \textbf{What it hands to the next stage} \\
    \hline
    1 & Formulate the question  & a question naming a group and a measurement \\
    2 & Collect or acquire data & the raw values, one per individual \\
    3 & Organize and represent  & a table that makes the pattern visible \\
    4 & Analyze and communicate & a conclusion in context --- and the next question \\
    \hline
    \end{tabular}
  \end{center}
  \vspace{8pt}
  \begin{center}
    {\large Each stage can use only what the stage before it \redink{handed over}.}\\[6pt]
    {\small\color{slate} If stage 1 never names the group, stage 4 cannot say who the
    conclusion is about.}
  \end{center}
\end{frame}

% ── ROW 2 — THE THREE TESTS ──────────────────────────────────────────────────
\begin{frame}[t, plain]
  \ceruleanheader{Stage 1 --- a question data can actually answer}
  \sectionlabel{notes row 2 $\cdot$ Three Tests $\cdot$ problems 5--8}
  \vspace{4pt}
  \begin{itemize}\setlength{\itemsep}{5pt}
    \item \textbf{1. Vary?} Would two individuals answer differently?
    \item \textbf{2. Group?} Who is the question about?
    \item \textbf{3. Record?} ``For each one I will write down \ldots'' --- finish that
          sentence. What you write down is the \textbf{\color{cerulean}variable of interest}.
  \end{itemize}
  \vspace{8pt}
  \begin{block}{Problem 8 --- test it before you touch it: ``How many minutes is a Riverbend student's trip to school?''}
    Vary? Yes. Group? Riverbend students. Record? Minutes. \redink{It passes all three.}
    Nothing to repair --- and a repair here would only change the subject.
  \end{block}
\end{frame}

% ── ROW 3 — FREQUENCY TABLE ──────────────────────────────────────────────────
\begin{frame}[t, plain]
  \ceruleanheader{Stages 2 and 3 --- collect, then organize}
  \sectionlabel{notes row 3 $\cdot$ Frequency Table $\cdot$ problems 9--11}
  \vspace{2pt}
  \begin{center}
    \small
    \renewcommand{\arraystretch}{1.25}
    \begin{tabular}{@{} l c l @{}}
    \hline
    \textbf{How they get to school} & \textbf{Frequency} & \textbf{Relative frequency} \\
    \hline
    Bus  & 27 & $27/60 = 0.45 = 45\%$ \\
    Car  & 15 & $15/60 = 0.25 = 25\%$ \\
    Walk & 12 & $12/60 = 0.20 = 20\%$ \\
    Bike &  6 & $6/60 = 0.10 = 10\%$ \\
    \hline
    \textbf{Total} & \textbf{60} & \textbf{100\%} \\
    \hline
    \end{tabular}
  \end{center}
  \vspace{6pt}
  \begin{center}
    {\color{cerulean}\bfseries What did stage 3 add? \redink{Nothing.} The counts were already
    there --- stage 3 made them readable.}\\[4pt]
    {\small Problem 11: $0.10 \times 900 = 90$ students. Now say it in a sentence that names
    the group.}
  \end{center}
\end{frame}

% ── ROW 4a — THE QUESTION DECIDES THE TYPE ───────────────────────────────────
\begin{frame}[t, plain]
  \ceruleanheader{The question decides the type --- before any data exist}
  \sectionlabel[redacc]{notes row 4 $\cdot$ The Question Decides $\cdot$ problems 12--13}
  \vspace{4pt}
  \begin{center}
    {\large The \emph{same} 30 students. Two questions. Two types.}
  \end{center}
  \vspace{6pt}
  \begin{columns}[t]
    \begin{column}{0.47\textwidth}
      \begin{block}{``How many mornings a week do you eat breakfast at school?''}
        Record a \textbf{count of mornings} $\to$ \textbf{\color{cerulean}quantitative}
        $\to$ summarize with a \textbf{mean}.
      \end{block}
    \end{column}
    \begin{column}{0.47\textwidth}
      \begin{block}{``Do you ever eat breakfast at school?''}
        Record \textbf{yes or no} $\to$ \textbf{\color{cerulean}categorical}
        $\to$ summarize with \textbf{percents}.
      \end{block}
    \end{column}
  \end{columns}
  \vspace{10pt}
  \begin{center}
    {\large You do \redink{not} have to collect anything first. What decided the type?
    \redink{The question.}}
  \end{center}
\end{frame}

% ── ROW 4b — THE CRUX: the hook settles here ─────────────────────────────────
\begin{frame}[t, plain]
  \ceruleanheader{Correct arithmetic. Unsupported claim.}
  \sectionlabel[redacc]{notes row 4 $\cdot$ problems 14--15 $\cdot$ the whole point of today}
  \vspace{4pt}
  \begin{itemize}\setlength{\itemsep}{4pt}
    \item \textbf{Problem 14 --- vote again: yes, no, or can't tell?}
    \item The council's numbers: $6/60 = 10\%$, about $90$ of $900$. \redink{All correct.}
    \item The question they asked was \emph{``How do you usually get to school?''} ---
          about \redink{now}. The word \emph{would} is nowhere in it.
  \end{itemize}
  \vspace{6pt}
  \begin{block}{Where did the problem enter? \redink{Stage 1.} (PS.DC.1b)}
    A conclusion answers \textbf{only the question that was asked}, about \textbf{only the
    group that was asked}. The bad sentence showed up at stage 4 --- but it \emph{entered} at
    stage 1, and no amount of correct arithmetic repairs that.
  \end{block}
  \vspace{4pt}
  {\small \textbf{Problem 15 --- the counterweight.} The cart's $14/40 = 35\%$ is also
  correct. ``35\% of Riverbend students buy a bagel every morning.'' That one entered at
  \redink{stage 2} --- only cart customers were ever asked, on one morning.}
\end{frame}

% ── WE DO — a LIVE surface: task and data only, no answers ───────────────────
\begin{frame}[t, plain]
  \ceruleanheader{Guided Practice: The Cycle Turns}
  \sectionlabel[goldacc]{We do $\cdot$ problems 16--19 $\cdot$ you hold the pen}
  The council goes back to \textbf{stage 1} with the question the first survey could not
  answer: \emph{``Would you ride a bike to school if there were racks?''} They ask
  \textbf{80} students; \textbf{24} say yes.
  \begin{itemize}\setlength{\itemsep}{5pt}
    \item[16.] The variable of interest and its type --- \emph{before} you look at the 80
          answers. How do you know?
    \item[17.] What percent said yes? About how many of the 900 is that?
    \item[18.] Write the conclusion the council can hand the district: the group asked, the
          number, what was measured.
    \item[19.] Suppose the 80 had all been standing at the bike rack. Which stage would the
          problem have entered at --- and what could the report no longer claim?
  \end{itemize}
  \begin{block}{The question behind the question}
    If you wrote \emph{stage 4} in 19, you named where the sentence appeared --- not where the
    problem entered.
  \end{block}
\end{frame}

% ── GUIDED PRACTICE problem 18 — shown AFTER the class has worked it ─────────
\begin{frame}[t, plain]
  \ceruleanheader{What the second survey lets them say --- and what it does not}
  \sectionlabel[goldacc]{Guided Practice, problem 18 $\cdot$ a conclusion is a sentence, not a number}
  \vspace{4pt}
  \begin{block}{The council's sentence}
    ``Of the \textbf{80} Riverbend students we asked, \textbf{30\%} said they would ride a
    bike to school if there were racks. If those 80 are typical, that is about \textbf{270} of
    our \textbf{900} students.''
  \end{block}
  \vspace{8pt}
  \begin{columns}[t]
    \begin{column}{0.47\textwidth}
      \begin{block}{``270 students will ride bikes.''}
        \redink{Not allowed.} 270 is a prediction from 80 answers --- nobody counted it, and
        nobody promised to ride.
      \end{block}
    \end{column}
    \begin{column}{0.47\textwidth}
      \begin{block}{``Riverbend families want bike racks.''}
        \redink{Not allowed.} Families were never asked, and racks were never the question.
      \end{block}
    \end{column}
  \end{columns}
\end{frame}

% ── DEBRIEF ──────────────────────────────────────────────────────────────────
\begin{frame}[t, plain]
  \ceruleanheader{Debrief: where did the problem enter?}
  \sectionlabel{10 minutes $\cdot$ pens down, papers up --- we say these aloud}
  \vspace{4pt}
  \begin{enumerate}\setlength{\itemsep}{5pt}
    \item The council's claim --- say why a correct $10\%$ still does not support it.
    \item Problem 12: nothing has been collected yet. How do you already know the type?
    \item Problem 15: name \emph{two} claims the cart data cannot support.
  \end{enumerate}
  \vspace{6pt}
  \begin{block}{Cold check --- hands down, thirty seconds}
    The yearbook staff surveyed only the students in the yearbook class, then wrote:
    \emph{Riverbend seniors prefer song lyrics in the yearbook.} \textbf{Which stage did the
    problem enter at, and why does that stop the conclusion?}
  \end{block}
\end{frame}

% ── YOU DO — the homework is the individual practice, started in class ───────
\begin{frame}[t, plain]
  \ceruleanheader{You do: start the homework}
  \sectionlabel[goldacc]{Last 10 minutes $\cdot$ on your own $\cdot$ finish it at home}
  \begin{itemize}\setlength{\itemsep}{5pt}
    \item \textbf{Lakeside Farmers Market} --- about 800 Saturday shoppers, 50 asked one
          Saturday.
    \item Name and classify the variable; the frequency table; scale 10\% to 800 and say it
          in a sentence.
    \item Item 4: the \emph{same} 50 shoppers, two questions --- two types. Decide before
          anything is collected.
    \item Item 6: every number in her report is correct and the conclusion still does not
          follow. Which stage did it enter at?
  \end{itemize}
  \begin{block}{Silent and alone --- I am circulating. Packet pages, scored out of 12, due at the start of the first class after two study halls.}
    Start with \textbf{1, 4, and 6}. Item 6 is the one I am reading over your shoulder --- if
    you wrote ``stage 4,'' flag me before you leave.
  \end{block}
\end{frame}

% ── CLOSE ────────────────────────────────────────────────────────────────────
{
\setbeamercolor{background canvas}{bg=cerulean}
\begin{frame}[plain]
  \vspace{1.5cm}\hspace{0.8cm}%
  \begin{minipage}{0.86\textwidth}
    {\color{goldacc}\bfseries\small BEFORE YOU GO}\\[0.7cm]
    {\color{white}\bfseries\Large The question decides what kind of data you get --- before any
    exist. And a percent can be perfectly correct and still not let you say the sentence you
    wanted to say.}\\[0.9cm]
    {\color{frostmid}\small Next: Lesson 1.2 --- the council asked 60 students and then talked
    about all 900. Those two groups have names, and so do the numbers computed from each.}
  \end{minipage}
\end{frame}
}

\end{document}
